\chapter{Conception}\label{chap:conception}

\section{Acteurs et besoins}
Les trois profils de l'atelier~2 sont inchangés (tableau~\ref{tab:acteurs}).

\begin{table}[htbp]
\centering
\caption{Acteurs et fonctionnalités}\label{tab:acteurs}
\begin{tabularx}{\linewidth}{@{}lX@{}}
\toprule
\textbf{Acteur} & \textbf{Fonctionnalités}\\
\midrule
Visiteur & Parcourir la vitrine, filtrer par catégorie, rechercher, ouvrir la fiche d'un produit, créer un compte, se connecter.\\
Client & Tout ce que fait le visiteur, plus : ajouter un produit au panier, modifier une quantité, retirer une ligne, vider le panier.\\
Administrateur & Créer, modifier, retirer de la vente et supprimer des produits ; gérer les catégories.\\
\bottomrule
\end{tabularx}
\end{table}

\section{Modèle du domaine}
Le modèle est identique à celui de l'atelier~2 : \code{Client} et \code{Administrateur} spécialisent la classe
abstraite \code{Internaute} ; un client possède un panier, composé de lignes qui référencent des produits ; la
vitrine présente des catégories qui regroupent des produits. La figure~\ref{fig:classes} le représente ; la
version Mermaid, lisible directement sur GitHub, figure en annexe (figure~\ref{fig:mermaid-metier}).

\begin{landscape}
\thispagestyle{empty}
\centering
\resizebox{\linewidth}{!}{\begin{tikzpicture}[
  % UML class box: title part (tinted) + attribute/method part
  cls/.style={
    rectangle split, rectangle split parts=2,
    draw=black!70, thick, fill=white,
    rectangle split part fill={teal!15, white},
    font=\footnotesize, text width=5.3cm, align=left,
    inner xsep=6pt, inner ysep=4pt},
  abs/.style={cls, draw=teal!60!black},
  mult/.style={font=\scriptsize, fill=white, inner sep=1pt},
  rel/.style={font=\scriptsize\itshape, fill=white, inner sep=1.5pt}
]

% --- Internaute hierarchy ---
\node[abs] (int) at (-2.9,0) {
  \textbf{\textit{Internaute}} \\ {\scriptsize\textit{<<abstract>>}}
  \nodepart{second}
  - id : Long\\
  - nom, prenom : String\\
  - email : String (unique)\\
  - motDePasseHash : String\\
  - dateInscription : LocalDateTime};

\node[cls] (adm) at (-6.0,-5.4) {
  \textbf{Administrateur}
  \nodepart{second}
  - matricule : String};

\node[cls] (cli) at (0.2,-5.4) {
  \textbf{Client}
  \nodepart{second}
  - adresse : String\\
  - telephone : String};

% --- Cart ---
\node[cls] (pan) at (8.8,-5.4) {
  \textbf{Panier}
  \nodepart{second}
  - id : Long\\
  - dateCreation : LocalDateTime\\
  - dateModification : LocalDateTime\\
  + ajouterProduit(Produit, int)\\
  + modifierQuantite(Produit, int)\\
  + retirerProduit(Produit)\\
  + vider()\\
  + getTotal() : BigDecimal};

\node[cls] (lig) at (17.4,-5.4) {
  \textbf{LignePanier}
  \nodepart{second}
  - id : Long\\
  - quantite : int\\
  - prixUnitaire : BigDecimal\\
  + getSousTotal() : BigDecimal};

% --- Catalogue ---
\node[cls] (pro) at (17.4,-11.6) {
  \textbf{Produit}
  \nodepart{second}
  - id : Long\\
  - reference : String (unique)\\
  - nom, description : String\\
  - prix : BigDecimal\\
  - stock : int\\
  - imageUrl : String\\
  - actif : boolean\\
  - version : long (@Version)};

\node[cls] (cat) at (8.8,-11.6) {
  \textbf{Categorie}
  \nodepart{second}
  - id : Long\\
  - nom : String\\
  - description : String};

\node[cls] (vit) at (0.2,-11.6) {
  \textbf{Vitrine}
  \nodepart{second}
  - id : Long\\
  - titre : String\\
  - description : String\\
  - active : boolean};

% --- Inheritance (hollow triangle) ---
\draw[-{Triangle[open,length=8pt,width=8pt]}, thick] (adm.north) -- (int.south);
\draw[-{Triangle[open,length=8pt,width=8pt]}, thick] (cli.north) -- (int.south);

% --- Client 1 - 1 Panier ---
\draw[thick] (cli.east) -- (pan.west)
  node[mult, pos=0.12, above] {1}
  node[mult, pos=0.88, above] {1}
  node[rel, midway, below] {poss\`ede};

% --- Panier composition LignePanier ---
\draw[{Diamond[length=9pt,width=6pt,fill=black]}-, thick] (pan.east) -- (lig.west)
  node[mult, pos=0.12, above] {1}
  node[mult, pos=0.88, above] {0..*}
  node[rel, midway, below] {contient};

% --- LignePanier -> Produit ---
\draw[-{Stealth[length=7pt]}, thick] (lig.south) -- (pro.north)
  node[mult, pos=0.12, right] {0..*}
  node[mult, pos=0.88, right] {1}
  node[rel, midway, right] {r\'ef\'erence};

% --- Categorie aggregation Produit ---
\draw[{Diamond[length=9pt,width=6pt,fill=white]}-, thick] (cat.east) -- (pro.west)
  node[mult, pos=0.12, above] {1}
  node[mult, pos=0.88, above] {0..*}
  node[rel, midway, below] {regroupe};

% --- Vitrine aggregation Categorie ---
\draw[{Diamond[length=9pt,width=6pt,fill=white]}-, thick] (vit.east) -- (cat.west)
  node[mult, pos=0.12, above] {1}
  node[mult, pos=0.88, above] {0..*}
  node[rel, midway, below] {pr\'esente};

% --- Administrateur manages Produit (dependency, dashed) ---
\coordinate (dep1) at (-6.0,-15.2);
\coordinate (dep2) at (17.4,-15.2);
\draw[-{Stealth[length=7pt]}, thick, dashed]
  (adm.south) -- (dep1) -- (dep2) -- (pro.south);
\node[rel] at (5.7,-15.2) {g\`ere};

\end{tikzpicture}
}
\captionof{figure}{Diagramme de classes du domaine (inchangé depuis l'atelier 2)}\label{fig:classes}
\end{landscape}

\section{Règles métier}
Les règles restent portées par l'entité \code{Panier} et par les contraintes de la base. Aucune n'est dupliquée dans
un bean : c'est la raison pour laquelle le remplacement de la couche web n'a demandé aucune modification du métier
(tableau~\ref{tab:regles}).

\begin{table}[htbp]
\centering
\caption{Règles métier et double protection}\label{tab:regles}
\begin{tabularx}{\linewidth}{@{}lXl@{}}
\toprule
\textbf{N°} & \textbf{Règle} & \textbf{Garde-fou en base}\\
\midrule
R1 & Une quantité de ligne est comprise entre 1 et 99. & \code{CHECK quantite >= 1}\\
R2 & Seul un produit actif peut être ajouté au panier. & ---\\
R3 & La quantité d'une ligne ne dépasse pas le stock disponible. & \code{CHECK stock >= 0}\\
R4 & Ajouter un produit déjà présent fusionne la ligne. & \code{UNIQUE (panier, produit)}\\
R5 & Le prix de la ligne est figé au moment de l'ajout. & \code{CHECK prix\_unitaire >= 0}\\
R6 & Le total est la somme des sous-totaux, en \code{BigDecimal}. & ---\\
R7 & L'adresse électronique est unique, en minuscules. & \code{UNIQUE (email)}\\
R8 & La référence d'un produit est unique, en majuscules. & \code{UNIQUE (reference)}\\
R9 & Un produit présent dans un panier ne peut pas être supprimé. & clé étrangère\\
R10 & Une catégorie qui contient des produits ne peut pas être supprimée. & clé étrangère\\
R11 & Un internaute a obligatoirement un rôle. & \code{NOT NULL} + \code{CHECK}\\
\bottomrule
\end{tabularx}
\end{table}

\section{Héritage}
La stratégie \code{SINGLE\_TABLE} est conservée : l'authentification lit un internaute par son adresse sans savoir
s'il est client ou administrateur, ce qui se résout en une requête sans jointure. Le détail de la comparaison avec
\code{JOINED} et \code{TABLE\_PER\_CLASS} figure dans le rapport de l'atelier~2.

\section{Architecture JSF}
Dans l'atelier~2, une servlet recevait la requête, appelait un service, plaçait des attributs dans la requête puis
transmettait à une JSP. Avec JSF, le contrôleur frontal est unique (\code{FacesServlet}) ; le rôle de chaque servlet
est repris par un \textbf{bean CDI}, lié à une page par des expressions EL. La figure~\ref{fig:beans} montre les trois
niveaux ; le tableau~\ref{tab:correspondance} donne la correspondance exacte avec l'atelier~2.

\begin{figure}[htbp]
\centering
\resizebox{\linewidth}{!}{% Pages -> beans -> services. Included inside a resizebox; no \begin{document} here.
\begin{tikzpicture}[
  node distance=0.9cm and 2.4cm,
  box/.style={draw=atlasteal,fill=codebg,rounded corners=2pt,minimum width=3.6cm,minimum height=0.85cm,align=center,font=\small},
  page/.style={box,draw=atlasink,fill=white},
  bean/.style={box,fill=atlasteal!12},
  svc/.style={box,fill=atlassaffron!18,draw=atlassaffron!80!black},
  arr/.style={-{Stealth[length=2mm]},thick,atlasink!70},
  col/.style={font=\bfseries\small\color{atlasteal}}]
  \node[page] (p1) {connexion.xhtml};
  \node[page,below=of p1] (p2) {inscription.xhtml};
  \node[page,below=of p2] (p3) {vitrine.xhtml\\detail.xhtml};
  \node[page,below=of p3] (p4) {panier.xhtml};
  \node[page,below=of p4] (p5) {admin/produits.xhtml\\admin/produit-form.xhtml};
  \node[page,below=of p5] (p6) {admin/categories.xhtml\\admin/categorie-form.xhtml};

  \node[bean,right=of p1] (b1) {ConnexionBean\\\scriptsize ViewScoped};
  \node[bean,right=of p2] (b2) {InscriptionBean\\\scriptsize ViewScoped};
  \node[bean,right=of p3] (b3) {VitrineBean, DetailBean\\\scriptsize ViewScoped};
  \node[bean,right=of p4] (b4) {PanierBean\\\scriptsize ViewScoped};
  \node[bean,right=of p5] (b5) {ProduitAdminBean\\\scriptsize ViewScoped};
  \node[bean,right=of p6] (b6) {CategorieAdminBean\\\scriptsize ViewScoped};

  \node[svc,right=of b1] (s1) {InternauteService};
  \node[svc] (s3) at ($(b3-|s1)$) {CatalogueService};
  \node[svc] (s4) at ($(b4-|s1)$) {PanierService};
  \node[svc] (s5) at ($(b5-|s1)$) {GestionCatalogueService};

  \draw[arr] (p1)--(b1); \draw[arr] (p2)--(b2); \draw[arr] (p3)--(b3);
  \draw[arr] (p4)--(b4); \draw[arr] (p5)--(b5); \draw[arr] (p6)--(b6);
  \draw[arr] (b1)--(s1); \draw[arr] (b2.east)--(s1.south west);
  \draw[arr] (b3)--(s3); \draw[arr] (b4)--(s4);
  \draw[arr] (b5)--(s5); \draw[arr] (b6.east)--(s5.south west);

  \node[col,above=0.3cm of p1] {Vues Facelets};
  \node[col,above=0.3cm of b1] {Beans CDI (@Named)};
  \node[col,above=0.3cm of s1] {Services (atelier 2)};
\end{tikzpicture}
}
\caption{Pages Facelets, beans CDI et services}\label{fig:beans}
\end{figure}

\begin{table}[htbp]
\centering
\caption{Correspondance servlets $\rightarrow$ beans}\label{tab:correspondance}
\begin{tabularx}{\linewidth}{@{}llX@{}}
\toprule
\textbf{Atelier 2} & \textbf{Atelier 3} & \textbf{Pages}\\
\midrule
\code{ClientServlet} & \code{ConnexionBean}, \code{InscriptionBean} & \code{connexion.xhtml}, \code{inscription.xhtml}\\
\code{VitrineServlet} & \code{VitrineBean}, \code{DetailBean} & \code{vitrine.xhtml}, \code{detail.xhtml}\\
\code{PanierServlet} & \code{PanierBean} & \code{panier.xhtml}\\
\code{ProduitServlet} & \code{ProduitAdminBean} & \code{admin/produits.xhtml}, \code{admin/produit-form.xhtml}\\
\code{CategorieServlet} & \code{CategorieAdminBean} & \code{admin/categories.xhtml}, \code{admin/categorie-form.xhtml}\\
\bottomrule
\end{tabularx}
\end{table}

Le sujet suggère un bean par gestion. Je l'ai suivi, avec un découpage plus fin que « un bean par entité » : un bean
est lié à un \textit{écran} et à son vocabulaire, pas à une table. Le panier, par exemple, manipule des lignes et des
produits mais n'a qu'un bean, parce qu'un seul écran le commande.

\fig{03-structure-projet}{Structure du projet dans Eclipse : beans, pages et composants}{fig:structure}
